\section{Categorical semantics of e-graphs}
\label{sec:categorical-semantics}

This section recalls the categorical semantics of e-graphs as morphisms in semilattice-enriched categories, due to~Tiurin et al.~\cite{tiurin2025equivalencehypergraphsdporewriting}, and then explains why we pass from monoidal categories to \emph{multicategories}.
Informally, a semilattice is a set equipped with an associative, commutative, and idempotent binary operation, written $\join$; a typical example is the set of finite non-empty subsets of a given set, with $\join \coloneqq \cup$.
The extensional content of an acyclic e-graph can be understood as a finite, non-empty collection of equivalent terms of a given type $A$.
For closed terms, such a collection is naturally represented by an element of a hom-semilattice $\enriched{C}(\I,A)$ in a cartesian semilattice-enriched category $\enriched{C}$; more generally, open terms in a context $\Gamma$ are represented in $\enriched{C}(\Gamma,A)$.
The join $M \join N$ records that $M$ and $N$ have been identified within the same e-class.
For example, the e-graph (b) of~\autoref{fig:e-graph-example} is the term
\begin{align}
\label{eq:e-graph-example-sigma-term}
(a \otimes ((2\fatsemi\comonoid) \otimes (1\fatsemi \id))\fatsemi\sym)\fatsemi((* \otimes \counit \join (\id \otimes \counit \otimes \id)\fatsemi<\!\!<) \otimes \id)\fatsemi/~,
\end{align}
where $\join$ encodes the identification of the subterms rooted at $*$ and $<\!\!<$.
As~\eqref{eq:e-graph-example-sigma-term} already shows, such raw terms are unwieldy: the arguments of an operation must be threaded through tensors, symmetries, and copy nodes before they reach that operation.
This term may also be contrasted with the term (b) in~\cref{lst:terms-e-graph}---the latter does not mention any $\id$ or $\fatsemi$ explicitly.
The type theory we develop in this paper makes this boilerplate implicit.

E-graphs are usually considered over an algebraic theory presented by a pair $(\Sigma, \mathcal{E})$, consisting of an algebraic signature $\Sigma$ and a set of equations $\mathcal{E}$.
An algebraic signature is a set of operations $c : n \to 1$ with arity $n$ and coarity $1$.
For example, the presentation for the algebraic theory used in~\cref{fig:e-graph-example} is given by
\begin{gather*}
\Sigma = \{\bar{n} : 0 \to 1 \mid n \in \mathbb{N}_{> 0}\} \cup \{ * : 2 \to 1, / : 2 \to 1, <\!\!< : 2 \to 1, a : 0 \to 1\},\\
\mathcal{E} = \{x * 2 = x <\!\!<1, (x * y) / z =  x * (y / z), x / x = 1, x * 1 = x\}.
\end{gather*}
The restriction to a single output is part of the conventional account of e-classes and e-nodes in traditional e-graphs.
Tiurin et al.~\cite{tiurin2025equivalencehypergraphsdporewriting} generalise from algebraic theories to symmetric monoidal theories, presented by signatures containing operations $c : n \to m$ with arbitrarily many outputs.
Equational reasoning in such theories is relevant to applications including probabilistic programming and quantum computing, where representing multi-output operations by ordinary tree-shaped terms is often impractical.

Algebraic and symmetric monoidal theories have clear categorical interpretations as Lawvere theories and PROPs, respectively.
Lawvere theories can be recovered as cartesian PROPs, so we describe the construction for PROPs.
This gives an interpretation of equivalence classes of terms as morphisms in a given PROP; ordinary hom-sets, however, have no operation representing the formation of an e-class.
To obtain one, we enrich the syntactic PROP generated by $\Sigma$ over the category of semilattices.
Then, the equations of $\mathcal{E}$ are implemented as rewrites on an enriched PROP.

\begin{definition}[Semilattice]
	A \emph{semilattice} is a set equipped with an operation $\join$ that is associative, commutative, and idempotent.
	A \emph{homomorphism} of semilattices is a map $h$ with $h(s \join s') = h(s) \join h(s')$.
	Semilattices and their homomorphisms form a category $\catname{SLat}$.
\end{definition}

\begin{proposition}[Special case of~{\cite[Proposition 6.4.6]{Borceux_1994}}]
\label{prop:free_slat}
	The forgetful functor $U \coloneqq \catname{SLat}(\I,-) : \catname{SLat} \to \catname{Set}$ has a strong monoidal left adjoint $F : \catname{Set} \to \catname{SLat}$, sending a set $A$ to the free semilattice on $A$.
\end{proposition}

This monoidal functor induces a 2-functor $\enriched{F} : \catname{Cat} \to \catname{SLat}\text{-}\catname{Cat}$.
It leaves the objects of each category $\mathcal{C}$ unchanged and applies $F$ to every hom-set, yielding an $\catname{SLat}$-category $\enriched{C}$.
Composition in $\enriched{C}$ is a semilattice homomorphism in each argument: whenever $f,f' : A \to B$ and $g,g' : B \to C$, we have
\begin{align}
(f \join f') \fatsemi g
  &= (f \fatsemi g) \join (f' \fatsemi g),
&
f \fatsemi (g \join g')
  &= (f \fatsemi g) \join (f \fatsemi g').
\label{eq:congruence_1}
\end{align}
As a 2-functor, free enrichment preserves adjunctions; because $F$ is strong monoidal, it also transports monoidal structure.
Thus, if $\mathcal{C}$ is monoidal, then $\enriched{C}$ is a monoidal $\catname{SLat}$-category, and tensoring is a semilattice homomorphism in each argument:
\begin{align}
(f \join f') \otimes g
  &= (f \otimes g) \join (f' \otimes g),
&
f \otimes (g \join g')
  &= (f \otimes g) \join (f \otimes g').
\label{eq:congruence_2}
\end{align}
Furthermore, every $\catname{SLat}$-functor $G : \enriched{C} \to \enriched{D}$ preserves joins:
\begin{equation}
G(f \join g) = G(f) \join G(g).
\label{eq:congruence_3}
\end{equation}
Equivalently, $\enriched{C}$ can be constructed directly from the signature $\Sigma$ by forming enriched $\Sigma$-terms---ordinary $\Sigma$-terms together with $\join$---and quotienting by the equations of a symmetric monoidal category and an $\catname{SLat}$-category.

In the e-graph interpretation, $M \join N$ records that $M$ and $N$ have been identified within one e-class.
Consequently, the multilinearity in~\cref{eq:congruence_1} is exactly the usual congruence property.
In conventional notation, if $x_1 \cong y_1, \ldots, x_n \cong y_n$, then $f(x_1, \ldots, x_n) \cong f(y_1, \ldots, y_n)$.
This is the relation maintained by traditional e-graphs~\cite{EggPaper}.
Equations~\eqref{eq:congruence_2} and~\eqref{eq:congruence_3} give the corresponding congruences for the additional categorical structure present in generalised e-graphs.
Whereas the e-graph literature often adds such congruences in an ad hoc manner~\cite{slotted-egraphs}, here they follow uniformly from enrichment over $\catname{SLat}$.
In particular,~\cref{eq:congruence_2} is the congruence law for the tensor product.
If $\mathcal{C}$ is closed monoidal, its free enrichment $\enriched{C}$ is closed monoidal as well: applying $F$ to the ordinary currying bijection produces an isomorphism of hom-semilattices.
For closed monoidal $\catname{SLat}$-categories,~\cref{eq:congruence_3} therefore also yields congruence for currying.
Such congruence is used by extensions of e-graphs with explicit variables and binding~\cite{slotted-egraphs,tiurin2026categoricalegraphslambdacalculi}.

In this work we generalise further from Lawvere theories and PROPs to multicategories and multicategorical presentations, in which morphisms with multiple inputs are first-class.
The enrichment machinery applies equally to multicategories.
Monoidal categories are recovered as \emph{representable} multicategories, while additional actions on inputs recover symmetry and cartesian structure.

\subsection{Multicategories}
\label{sec:multicategories}

Multicategories were introduced by~Lambek~\cite{LambekMulticategories} as a generalisation of categories in which morphisms may have several inputs.

\begin{definition}[Multicategory]
\label{def:multicategory}
    A \emph{multicategory} $\mathcal{M}$ consists of the following data:
    \begin{enumerate}
        \item a collection of objects $\obj{\mathcal{M}}$;
        \item for each $n \geq 0$ and objects $A_1, \ldots, A_n, B$, a set of multimorphisms
        \[\mathcal{M}(A_1, \ldots, A_n; B)\,;\]
        \item for each object $A$, an identity morphism $\id_A \in \mathcal{M}(A; A)$;
        \item for any object $C$ and lists of objects $(B_1, \ldots, B_m)$ and $\vec A_i=(A_{i1}, \ldots, A_{in_i})$ for $1 \leq i \leq m$, a composition operation
        \begin{align*}
            \mathcal{M}(B_1, \ldots, B_m;C) \times \prod_{i=1}^{m}\mathcal{M}(\vec A_i; B_i)
            &\to \mathcal{M}(\vec A_1,\ldots,\vec A_m;C)\\
            (g, (f_1, \ldots, f_m)) &\mapsto g \circ (f_1, \ldots, f_m)
        \end{align*}
    \end{enumerate}
    subject to the following conditions:
    \begin{itemize}
        \item for any $f \in \mathcal{M}(A_1, \ldots, A_n; B)$,
        \[
        \id_{B} \circ (f) = f \qquad f \circ (\id_{A_{1}}, \ldots, \id_{A_n}) = f\,;
        \]
        \item for any $h,g_i,f_{ij}$,
        \[
        \begin{aligned}
        &(h \circ (g_1, \ldots, g_m)) \circ (f_{11}, \ldots, f_{mn_m})\\
        &\qquad={}
        h \circ (g_1 \circ (f_{11}, \ldots, f_{1n_1}), \ldots,
                  g_m \circ (f_{m1}, \ldots, f_{mn_m}))\,.
        \end{aligned}
        \]
    \end{itemize}
\end{definition}

A typical example is the multicategory of vector spaces and multilinear maps.
Multifunctors preserve identities and composition, while multinatural transformations satisfy the corresponding many-input naturality condition.
Together these form a $2$-category $\catname{Multicat}$; we spell out the definitions in~\Cref{app:multicategories}.
\begin{remark}
\label{remark:placed_composition}
    Equivalently, one can define a multicategory using the data (1)--(3) above, together with a placed binary composition
    \begin{gather*}
    \mathcal{M}(B_1, \ldots, B_m;C) \times \mathcal{M}(A_1, \ldots, A_n; B_i)\\
    \longrightarrow
    \mathcal{M}(B_1, \ldots, B_{i-1}, A_1, \ldots, A_n, B_{i+1}, \ldots, B_m;C),
    \qquad (g,f) \longmapsto g \circ_i f,
    \end{gather*}
    subject to the following conditions:
    \begin{itemize}
        \item $\id_{B} \circ_1 f = f$;
        \item $f \circ_{i} \id_{B_{i}} = f$;
        \item if $h$ is $n$-ary, $g$ is $m$-ary and $f$ is $k$-ary, then
        \[
        (h \circ_{i} g) \circ_{j} f = \begin{cases}
        (h \circ_{j} f) \circ_{i+k-1} g & j < i\\
        h \circ_{i} (g \circ_{j-i+1} f) & i \leq j < i + m\\
        (h \circ_{j-m+1} f) \circ_{i} g & j \geq i + m
        \end{cases}
        \]
    \end{itemize}
    The middle case of the last equation is the associativity law; the other two cases are the interchange laws.
    Here $1 \leq i \leq n$ and $1 \leq j \leq n+m-1$.
    If $j<i$, the second composition inserts $f$ into an input of $h$ that precedes the input occupied by $g$; if $j\geq i+m$, it inserts $f$ into an input following $g$.
    These two independent insertions may be performed in either order, giving the interchange laws.
    If $i\leq j<i+m$, the second composition instead inserts $f$ into one of the inputs of $g$, giving associativity.
    The three cases are depicted in~\Cref{fig:composition_multicat}.

    We will show that substitution in our type theory satisfies these laws and therefore defines multicategorical composition.
\end{remark}
\begin{figure}[t]
    \centering
    \adjustbox{width=0.4\linewidth}{
    \tikzfig{./figures/composition_assoc}
    }
    \caption{Placed composition in a multicategory.}
    \label{fig:composition_multicat}
\end{figure}
\begin{definition}[Tensor, representability]
\label{def:multicat_tensor}
    Given objects $A_1, \ldots, A_n$, a \emph{tensor product} of them in a multicategory $\mathcal{M}$ is an object $\bigotimes_{i}A_{i}$ together with a universal morphism $\xi \in \mathcal{M}(A_1, \ldots, A_{n}; \bigotimes_{i}A_{i})$ such that composition with $\xi$ induces bijections, natural in $D$ and multinatural in the $B_i$ and $C_j$:
    \[
    \mathcal{M}(B_{1}, \ldots, B_{m}, \bigotimes_{i}A_{i}, C_{1}, \ldots, C_{k}; D)
    \xrightarrow{-\circ_{m+1}\xi}
    \mathcal{M}(B_1, \ldots, B_m, A_{1}, \ldots, A_{n}, C_{1}, \ldots, C_{k}; D)\,.
    \]
    A unit $\I$ is a nullary tensor product.
    $\mathcal{M}$ is \emph{representable} if it has a chosen unit and a binary tensor for every pair of objects.
\end{definition}

\begin{proposition}[{\cite[Thm. 9.8]{hermida_representable_2000}}]
\label{prop:representable_multicategories}
    There is a 2-equivalence $\catname{RepMultCat} \simeq \catname{MonCat}$ between representable multicategories and (non-strict) monoidal categories with strong monoidal functors.
\end{proposition}

More explicitly, if $\mathcal{C}$ is the monoidal category corresponding to a representable multicategory $\mathcal{M}$, then its unary hom-sets are
$\mathcal{C}(A,B)=\mathcal{M}(A;B)$, and representability gives natural bijections
\[
\mathcal{M}(A_1,\ldots,A_n;B)
\cong
\mathcal{C}(A_1\otimes\cdots\otimes A_n,B).
\]


This correspondence extends to symmetric and cartesian monoidal categories by
equipping the multicategory with a compatible action on its inputs.
Concretely, a chosen subcategory $\mathfrak{S}$ of finite cardinals and functions acts by reindexing: every $a:m\to n$ induces a map
\[
a^* : \mathcal{M}(A_{a(1)},\ldots,A_{a(m)};B)
\longrightarrow \mathcal{M}(A_1,\ldots,A_n;B).
\]
These maps preserve identities, compose according to the composition in $\mathfrak{S}$, and commute with multicategorical composition.
If $\mathfrak{S}$ consists of bijections, this action merely permutes inputs.
If it contains all functions, non-injective maps allow an input to be duplicated, while non-surjective maps allow unused inputs to be discarded.
Thus the latter action supplies exchange, contraction, and weakening.

\begin{example}
\label{ex:symmetric_S_multicategory}
    If $\mathfrak{S}$ consists of all bijections, representable $\mathfrak{S}$-multicategories are $2$-equivalent to symmetric monoidal categories.
    The symmetry $\sym_{A,B} : A \otimes B \to B \otimes A$ is the unique morphism with $\sym_{A,B} \circ \xi_{A,B} = a^*(\xi_{B,A})$ for the transposition $a : \underline{2} \to \underline{2}$.
    The verification that $\sym$ is involutive and natural is given in~\Cref{app:multicategories}.
    Taking $\mathfrak{S}$ to be all functions instead yields cartesian categories---the case of conventional e-graphs.
\end{example}

Enriching multicategories proceeds as for categories: one replaces each hom-set by a hom-object in the enriching monoidal category and requires composition and structural actions to be morphisms of that category.
We follow the formulation of Johnson and Yau~\cite{johnson2023homotopytheoryenrichedmackey}, recalled in~\Cref{app:multicategories}.
For $\catname{SLat}$-enrichment, the hom-object $\enriched{M}(A_1,\ldots,A_n;B)$ is a semilattice, and the corresponding ordinary hom-set is its underlying set
\[
U\enriched{M}(A_1,\ldots,A_n;B)
=\catname{SLat}(\I,\enriched{M}(A_1,\ldots,A_n;B)).
\]
Thus every ordinary multimorphism is an element of the corresponding hom-semilattice, and every structural reindexing map $a^*$ preserves joins.
Composition is likewise \emph{multilinear}: it preserves joins in each argument separately.
In particular,
\[
 (f \join g) \circ (h_1,\ldots,h_n)
 = \bigl(f \circ (h_1,\ldots,h_n)\bigr)
   \join \bigl(g \circ (h_1,\ldots,h_n)\bigr),
\]
\[
f \circ (h_1, \ldots, g \join g', \ldots, h_n)
= \bigl(f \circ (h_1, \ldots, g, \ldots, h_n)\bigr)
  \join \bigl(f \circ (h_1, \ldots, g', \ldots, h_n)\bigr),
\]
or, using placed composition $\circ_i$,
\[
f_1 \circ_i (g_1 \join g_2)
= (f_1 \circ_i g_1) \join (f_1 \circ_i g_2),
\qquad
(f_1 \join f_2) \circ_i g
= (f_1 \circ_i g) \join (f_2 \circ_i g).
\]
For a representable $\catname{SLat}$-multicategory, the tensor product of morphisms $f \in \mathcal{M}(A; B)$ and $g \in \mathcal{M}(C; D)$ is the unique morphism $f \otimes g \in \mathcal{M}(A \otimes C; B \otimes D)$ such that
\[
(f \otimes g) \circ_1 \xi_{A,C} = \xi_{B,D} \circ (f,g).
\]
For $f_1,f_2:A\to B$ and $g:C\to D$, multilinearity gives
\begin{align*}
((f_1 \join f_2) \otimes g) \circ_1 \xi_{A,C}
 &= \xi_{B,D} \circ (f_1 \join f_2,g)\\
 &= \bigl(\xi_{B,D} \circ (f_1,g)\bigr)
      \join \bigl(\xi_{B,D} \circ (f_2,g)\bigr)\\
 &= \bigl((f_1\otimes g)\join(f_2\otimes g)\bigr)\circ_1\xi_{A,C}.
\end{align*}
The universal bijection is injective, and hence
\[
(f_1\join f_2)\otimes g=(f_1\otimes g)\join(f_2\otimes g).
\]

Finally, a \emph{closed} multicategory may be given by a similar universal property.
% A multicategory is \emph{closed} if for every $A,B$ it has an object
% $A\multimap B$ and an evaluation morphism
% \[
% \mathsf{ev}_{A,B}\in\mathcal M(A\multimap B,A;B)
% \]
% such that evaluation induces natural bijections
% \[
% \Lambda_A:\mathcal M(\Gamma,A;B)
% \xrightarrow{\cong}\mathcal M(\Gamma;A\multimap B).
% \]
% For an $\catname{SLat}$-multicategory these are required to be isomorphisms of
% semilattices.  Given $f:B\to C$, the covariant action of the internal hom is
% therefore
% \[
% A\multimap f
% \coloneqq
% \Lambda_A\bigl(f\circ_1\mathsf{ev}_{A,B}\bigr)
% :A\multimap B\longrightarrow A\multimap C.
% \]
% Hence, for $f,g:B\to C$, multilinearity of composition and preservation of
% joins by $\Lambda_A$ give
% \begin{align*}
% A\multimap(f\join g)
%  &=\Lambda_A\bigl((f\join g)\circ_1\mathsf{ev}_{A,B}\bigr)\\
%  &=\Lambda_A\bigl((f\circ_1\mathsf{ev}_{A,B})
%        \join(g\circ_1\mathsf{ev}_{A,B})\bigr)\\
%  &=(A\multimap f)\join(A\multimap g).
% \end{align*}
See~\cite{shulman2016categorical} for the full multicategorical treatment of
closedness.
We now have the machinery needed to define the type theory for
$\catname{SLat}$-multicategories.
